\documentclass[11pt,a4paper]{article}
\usepackage[ngerman]{babel}
\usepackage[TS1,T1]{fontenc}
\usepackage[utf8x]{inputenc}
\selectlanguage{ngerman}
\usepackage{theorem}
\usepackage[scaled=0.9]{helvet}
\usepackage{amsmath}
\usepackage{amssymb}
\usepackage{hyperref}
\usepackage{stmaryrd}
\usepackage{pgf,tikz}
\usepackage{relsize}
\usepackage{enumitem}
\usepackage{graphicx}
\usepackage{algpseudocode,amsmath,xifthen}

\newcounter{numb}
\theoremstyle{break}
\theorembodyfont{\upshape}
   	\newtheorem{exercise}{Exercise}[numb]
\setlength{\oddsidemargin}{0cm}
\setlength{\textwidth}{16cm}
\setlength{\textheight}{23cm}
\setlength{\topmargin}{-2cm}

\usetikzlibrary{shapes,arrows,automata,positioning,decorations.fractals}
\renewcommand\familydefault{\sfdefault}

\newcommand{\header}[2]{
\begin{minipage}[b]{\textwidth}
		\parbox[t]{5cm}{%
			\includegraphics[width=4cm]{unilogo}
			\hfill
		}
		\parbox[b]{11cm}{%
			%\scshape%
			Heinrich-Heine-University D\"usseldorf\\
			Computer Science Department\\
			Software Engineering and Programming Languages\\
			%Professor Dr.\ M.\ Leuschel
		Philipp K\"orner
}
		
\end{minipage}
	\begin{center}
		\bf
		Functional Programming -- WT 2024 / 2025\\
        Reading Guide {\thenumb}: {#1}
	\end{center}

    \pagestyle{empty}
    \paragraph{Timeline:} This unit should be completed by {#2}.
}

\setcounter{numb}{0}

\begin{document}
	
\begin{minipage}[b]{\textwidth}
		\parbox[t]{5cm}{%
			\includegraphics[width=4cm]{unilogo}
			\hfill
		}
		\parbox[b]{11cm}{%
			%\scshape%
			Heinrich-Heine-University D\"usseldorf\\
			Computer Science Department\\
			Software Engineering and Programming Languages\\
			%Professor Dr.\ M.\ Leuschel
		Philipp K\"orner
}
		
\end{minipage}
	\begin{center}
		\bf
		Functional Programming -- WISE 2020 / 2021\\
		Reading Guide 0: Overview
	\end{center}
	
	\pagestyle{empty}
	
	\section*{Note}
	The learning objectives indicate what is expected of you in the exam.
	In principle it does not matter which material you consult to achieve the learning objective.
	The deadline allows you to meaningfully participate in the exercises.
	
	Solving the exercises is voluntary.
	The exercises in this guide present tools and as such are elementary for subsequent course material.
	In addition these tools are an immense help for solving the exercises.
	
	You are free to proceed to the next unit, as soon as you complete this one.  
	
	\paragraph{Deadline:} This unit should be completed by 05.11.2020.
	
	\section{Material} 
	
	\begin{itemize}
		\item The Joy of Clojure, Chapter 1 - Clojure Philosophy
		\footnote{Do not be discouraged if you do not understand everything from this chapter immediately. The book is not especially suited for beginners, but gives a comprehensive overview of the topics we will focus on.}
		\item 00-syntax.clj
		\item Leiningen \url{https://github.com/technomancy/leiningen}
		\item Clojure Cheatsheet \url{https://clojure.org/api/cheatsheet}
		\item Clojure Newbie Guide \url{http://www.clojurenewbieguide.com/}
	\end{itemize}
	
	
	\section{Learning Objectives}
	
	After completing this unit you should be able to
	
	\begin{itemize}
		\item name the characteristics of functional programming.
		\item describe the downsides of object-oriented programming.
		\item use an editor with a REPL-plugin and paredit.
		\item list some Clojure-functions and explain their functionality.
	\end{itemize}
	
	\section{Exercises}
	
	
	\begin{exercise}[Leiningen]
		Install Leiningen (\url{https://github.com/technomancy/leiningen}).
		Please use the installation instructions provided on the linked website
		to ensure you are using the latest version.
		
		Find out how to generate a new project
		and familiarise yourself with the directory structure.
		Start a REPL-session as well.
	\end{exercise}
	
	
	\begin{exercise}[Paredit]
		Paredit is an important aid in programming in a LISP.
		A paredit-plugin exists for most editors.
		Make sure that your editor of choice has such a plugin
		and install or activate it.
		
		
		Paredit helps keep parentheses balanced.
		In order to work efficiently you should
		become acquainted with the keyboard shortcuts that allow for the following manipulations
		(copy \& paste or retyping are prohibited!):
		
		\begin{enumerate}[label=\alph*)]
			\item
			\texttt{(foo (bar) baz)} $\longrightarrow$ \texttt{(foo (bar baz))} \textit{(Right Slurp)}
			\item
			\texttt{(foo (bar baz))} $\longrightarrow$ \texttt{(foo bar (baz))} \textit{(Left Barf)}
			\item
			\texttt{(foo bar (baz))} $\longrightarrow$ \texttt{(foo (bar baz))} \textit{(Left Slurp)}
			\item
			\texttt{(foo (bar baz))} $\longrightarrow$ \texttt{(foo (bar) baz)} \textit{(Right Barf)}
			\item
			\texttt{(foo) (bar)} $\longrightarrow$ \texttt{(foo bar)} \textit{(Join)}
			\item
			\texttt{(foo bar)} $\longrightarrow$ \texttt{(foo) (bar)} \textit{(Split)}
			\item
			\texttt{foo} $\longrightarrow$ \texttt{(foo)} \textit{(Wrap)}
			\item
			\texttt{(foo (bar baz))} $\longrightarrow$ \texttt{(foo bar baz)} \textit{(Splice)}
			\item
			\texttt{(foo (bar baz))} $\longrightarrow$ \texttt{(foo)} \textit{(Delete / Killing)}
		\end{enumerate}
		
	\end{exercise}
	
	\begin{exercise}[REPL-Plugin]
		Many editors provide plugins
		to connect to open REPL sessions
		and execute code within their context.
		
		Find out if such a plugin exists for your chosen editor,
		start a REPL session and evaluate the following
		code (You do not have to understand the code at this point!)
		
		\begin{verbatim}
			(def fib-seq
			    (lazy-cat [0 1] (map + (rest fib-seq) fib-seq)))
			
			(take 10 fib-seq)
		\end{verbatim}
	\end{exercise}
	
	
	\section*{Questions}
	If you have any questions, please contact Philipp K"orner (\texttt{p.koerner@hhu.de}).
\end{document}

